\Askhsh{27031}{4}
\begin{ylh}{1}
    \item 2.1 Η έννοια της παραγώγου
    \item 2.4 Ρυθμός μεταβολής
    \item 3.7 Εμβαδόν επιπέδου χωρίου.
\end{ylh}
\begin{wrapfigure}[24]{r}{7cm} % Adjust line wrap number and width as needed
    \centering
    \begin{tikzpicture}
    \begin{axis}[
        axis lines = middle,
        xlabel = {$x$},
        ylabel = {$y$},
        xmin = -4.5, xmax = 0.5,
        ymin = -1, ymax = 10,
        xtick = {-4,-3,-2,-1,0},
        ytick = {0,1,2,3,4,5,6,7,8,9,10},
        grid = none, % No grid in the original image
        axis line style = {->}, % Arrows at the end of axes
        xlabel style={at={(current axis.right of origin)},anchor=north east},
        ylabel style={at={(current axis.above origin)},anchor=south east},
        tick label style = {font=\scriptsize},
        declare function={f(\x) = -1/3*\x^3;}, % Define f(x)
    ]
        % Plot the curve f(x) = -1/3 x^3
        \addplot[very thick, maincolor, domain=-4:0, samples=100] {f(x)};
        \node[maincolor, above right, xshift=5pt, yshift=5pt] at (axis cs:0,0) {$y = -\frac{1}{3}x^3$};

        % Define a specific alpha for visualization, e.g., alpha = -3
        \pgfmathsetmacro{\alphaVal}{-3}
        \pgfmathsetmacro{\yA}{-1/3 * (\alphaVal)^3} % f(\alphaVal) = -1/3 * (-3)^3 = 9
        \pgfmathsetmacro{\xM}{2/3 * \alphaVal} % M's x-coordinate is 2\alpha/3 = 2/3 * (-3) = -2

        % Point A
        \coordinate (A_point) at (axis cs:\alphaVal, \yA); % A(-3, 9)
        \node[above left, xshift=-2pt, yshift=-2pt] at (A_point) {$A\left(\alpha, -\frac{\alpha^3}{3}\right)$};
        \fill (A_point) circle (1.5pt);

        % Tangent line from A to M
        % Calculate tangent slope at alpha: f'(\alpha) = -\alpha^2 = -(-3)^2 = -9
        % Tangent line equation: y - yA = slopeA * (x - xA)
        % For plotting: y = slopeA * (x - xA) + yA
        \pgfmathsetmacro{\slopeA}{-(\alphaVal)^2}
        \addplot[thick, secondarycolor, domain=\xM:\alphaVal, samples=2] {\slopeA*(x-\alphaVal) + \yA}; % Plot tangent segment
        
        % Vertical dashed line from A to B (projection on x-axis)
        \coordinate (B_point) at (axis cs:\alphaVal, 0); % B(-3, 0)
        \draw[dashed] (A_point) -- (B_point);
        \fill (B_point) circle (1.5pt);
        \node[below, yshift=-2pt] at (B_point) {$B$};

        % Point M (x-intercept of the tangent)
        \coordinate (M_point) at (axis cs:\xM, 0); % M(-2, 0)
        \fill (M_point) circle (1.5pt);
        \node[below, yshift=-2pt] at (M_point) {$M$};

        % Origin O
        \coordinate (O_point) at (axis cs:0,0);
        \fill (O_point) circle (1.5pt);
        \node[below right, xshift=2pt, yshift=-2pt] at (O_point) {$O$};
        
        % "Ακτή" label below the x-axis
        \node[below, yshift=-5pt] at (axis cs:-1,0) {Ακτή};

        % Arrows representing the light beam, as shown in the original image.
        % Arrow coming down to A
        \draw[thick, ->, maincolor] (axis cs:\alphaVal, \yA+1.5) -- (A_point);
        % Arrow along the tangent (light beam originating from A, pointing towards M)
        \draw[thick, ->, maincolor] (A_point) -- (M_point);

    \end{axis}
    \end{tikzpicture}
\end{wrapfigure}

Δίνεται η συνάρτηση $f(x) = -\frac{1}{3}x^3$, με $x \in (-\infty, 0]$ και τυχαίο σημείο $A\left(\alpha, -\frac{\alpha^3}{3}\right)$ με $\alpha < 0$ της γραφικής της παράστασης.

\begin{alist}
\item Να βρείτε την εξίσωση της εφαπτομένης της $C_f$ στο σημείο $A$. \monades{06}
\item
    \begin{rlist}
    \item Ένα περιπολικό $A$ κινείται κατά μήκος της καμπύλης $y = -\frac{1}{3}x^3$, $x \le 0$ πλησιάζοντας την ακτή και ο προβολέας του φωτίζει κατευθείαν εμπρός (όπως φαίνεται στο σχήμα). Αν ο ρυθμός μεταβολής της τετμημένης του περιπολικού δίνεται από τον τύπο $\alpha'(t) = -\alpha(t)$, να βρείτε το ρυθμό μεταβολής της τετμημένης του σημείου $M$ της ακτής, στο οποίο πέφτουν τα φώτα του προβολέα τη χρονική στιγμή $t_0$, κατά την οποία το περιπολικό έχει τετμημένη $-3$. \monades{08}
    \item Να ερμηνεύσετε το πρόσημο του ρυθμού μεταβολής της τετμημένης του σημείου $M$. \monades{02}
    \end{rlist}
\item Να βρείτε το εμβαδόν του χωρίου $\Omega$, που περικλείεται από την γραφική παράσταση της συνάρτησης $f$, τον άξονα $x'x$ και την εφαπτομένη της $C_f$ στο σημείο της με τετμημένη $-3$. \monades{09}
\end{alist}